\section{Principal congruence subgroups of the Atkin--Lehner groups}
\label{sec:principal}

\subsection{Definitions}

\begin{definition}[Eichler orders]\label{def:eichler-order}
For $N \geq 1$ the \emph{Eichler order of level $N$} is
$\EichlerOrder{N} := \{\smallMat{a}{b}{Nc}{d} : a,b,c,d \in \ZZ\} \subseteq
\MatRing{2}{\ZZ}$, and $\AdelicEichlerOrder{N} := \EichlerOrder{N}
\otimes\ProfiniteIntegers \subseteq \MatRing{2}{\ProfiniteIntegers}$. For a prime
$p$ the \emph{Iwahori order} is $\IwahoriOrder{p} = \EichlerOrder{p}\otimes
\pAdicIntegers{p}$ (as in \cref{lem:local-fixed-points}).
\end{definition}

\begin{definition}[Principal congruence subgroups of level $(N;M)$]
\label{def:principal-NM}
Let $N, M \geq 1$ with $N$ squarefree. Define
\begin{align*}
  \GammaPrincipal[ad]{N}{M} &:= \AdelicProjection\{g \in \unit{\AdelicEichlerOrder{N}} :
  g \equiv 1 \bmod M\AdelicEichlerOrder{N}\} \subseteq \AdelicGroup,\\
  \GammaPrincipal[lift]{N}{M} &:= \GammaFull[lift]{M}\cap\GammaZero[lift]{NM}
  = \{\gamma \in \SL{2}{\ZZ} : \gamma \equiv 1 \bmod M\EichlerOrder{N}\}.
\end{align*}
\end{definition}

\begin{lemma}[{Lifted version of $\GammaPrincipal[ad]{N}{M}$}]
\label{lem:principal-lifted}
Let $N, M \geq 1$ with $N$ squarefree. The lifted version of
$\GammaPrincipal[ad]{N}{M}$ in the sense of \cref{def:versions} is
$\GammaPrincipal[lift]{N}{M}$. Moreover $\GammaPrincipal[ad]{N}{1} =
\GammaZero[ad]{N}$, $\GammaPrincipal[ad]{1}{M} = \GammaFull[ad]{M}$, and
$M\EichlerOrder{N}$ is the set of matrices $X \in \MatRing{2}{\ZZ}$ with
$X \equiv 0 \bmod M$ and $X_{21} \equiv 0 \bmod NM$.
\end{lemma}

\begin{proof}
The rational elements are described by \cref{lem:rational-points-glzhat};
the rest is immediate from the definitions.
\end{proof}

\begin{lemma}[$M\EichlerOrder{N}$ is a two-sided ideal]\label{lem:eichler-ideal}
Let $N, M \geq 1$. Then $M\EichlerOrder{N}$ is a two-sided ideal of
$\EichlerOrder{N}$, and for every prime $p$ the Iwahori order satisfies
$\IwahoriOrder{p} w = w\IwahoriOrder{p} =: \IwahoriRadical{p}$ and
$\IwahoriRadical{p}^{2} = p\IwahoriOrder{p}$, where $w = \smallMat{0}{1}{p}{0}$
and $\IwahoriRadical{p} = \{X \in \MatRing{2}{\pAdicIntegers{p}} :
p \mid X_{11}, X_{21}, X_{22}\}$ is the Jacobson radical of
$\IwahoriOrder{p}$. In particular $w$ normalises $\IwahoriOrder{p}$ and
$p^{k}\IwahoriOrder{p}$ for all $k \geq 0$.
\end{lemma}

\begin{proof}
The first statement is clear. One computes
$\smallMat{a}{b}{pc}{d}w = \smallMat{pb}{a}{pd}{pc}$ and
$w\smallMat{a}{b}{pc}{d} = \smallMat{pc}{d}{pa}{pb}$, both of which run
through $\IwahoriRadical{p}$. Then $\IwahoriRadical{p}^{2} =
\IwahoriOrder{p}w^{2}\IwahoriOrder{p} = p\IwahoriOrder{p}$ since $w^{2} = p$.
The quotient $\IwahoriOrder{p}/\IwahoriRadical{p} \cong \FF{p}\times\FF{p}$
(diagonal entries) is semisimple and $\IwahoriRadical{p}$ is topologically
nilpotent, so it is the Jacobson radical.
\end{proof}

\subsection{Normality and the core}

\begin{proposition}[{$\GammaPrincipal[ad]{N}{M}$ is normal in the
Atkin--Lehner group}]\label{prop:principal-normal}
\slogan{$\GammaPrincipal{N}{M}$ is normal in $\GammaZeroStar{N}$.}
Let $N, M \geq 1$ with $N$ squarefree. Then $\GammaPrincipal[ad]{N}{M}$ is a
normal compact open subgroup of $\MaxCompact{N} = \GammaZeroStar[ad]{N}$ of
trivial type, and $\GammaPrincipal[lift]{N}{M}$ is a normal subgroup of
$\GammaZeroStar[lift]{N}$.
\end{proposition}

\begin{proof}
$\GammaPrincipal[ad]{N}{M}$ is the image of the product over $p$ of the
groups $1 + p^{\val{p}{M}}(\EichlerOrder{N}\otimes\pAdicIntegers{p})$ (unit
groups), which is compact open. It is normalised by
$\unit{(\EichlerOrder{N}\otimes\pAdicIntegers{p})}$ because
$M\EichlerOrder{N}$ is a two-sided ideal, and by the local flips at
$p \mid N$ because they normalise $\IwahoriOrder{p}$ and its ideals
$p^{k}\IwahoriOrder{p}$ (\cref{lem:eichler-ideal}; the local flip differs
from $w$ by an element of $\unit{\IwahoriOrder{p}}$ up to scalars). By the
structure of the Atkin--Lehner group (\cref{lem:atkin-lehner-group}) these
generate $\MaxCompact{N}$ together with the centre. The lifted statement
follows by intersecting with $\RationalLiftGroup$
(\cref{lem:principal-lifted}).
\end{proof}

\begin{proposition}[$\GammaPrincipal{N}{M}$ is a core]\label{prop:principal-core}
\slogan{$\GammaPrincipal{N}{M}$ is the largest subgroup of $\GammaFull{M}$
normal in $\GammaZeroStar{N}$.}
Let $N, M \geq 1$ with $N$ squarefree. Then
\[
  \GammaPrincipal[ad]{N}{M} = \Core{\MaxCompact{N}}{\GammaFull[ad]{M}\cap\MaxCompact{N}},
  \qquad
  \GammaPrincipal[lift]{N}{M} = \Core{\GammaZeroStar[lift]{N}}{\GammaFull[lift]{M}}.
\]
\end{proposition}

\begin{proof}
The inclusions $\subseteq$ hold because the left hand sides are normal
(\cref{prop:principal-normal}) and contained in the principal congruence
subgroup. Conversely the core is contained in
$\GammaFull[ad]{M} \cap w\GammaFull[ad]{M}w^{-1}$ for every local flip
$w = \ALflip{p,N}$; at $p$ one has
$\MatRing{2}{\pAdicIntegers{p}} \cap w\MatRing{2}{\pAdicIntegers{p}}w^{-1} =
\IwahoriOrder{p}$, since conjugation by $w$ maps
$\smallMat{a}{b}{c}{d}$ to $\smallMat{d}{-c/p}{-pb}{a}$. Hence the local
component at $p$ of an element of the core lies in
$1 + p^{\val{p}{M}}\IwahoriOrder{p}$ (for $\val{p}{M} = 0$: in
$\unit{\IwahoriOrder{p}}$). For the lifted statement use that
$\ALmatrix{N}$ maps to $\GammaZero[ad]{N}\prod_{p\mid N}\ALflip{p,N}$
(\cref{lem:global-al}) and that $\GammaZero[ad]{N}$ normalises
$\GammaFull[ad]{M}$.
\end{proof}

\begin{lemma}[Strong approximation for Eichler orders]
\label{lem:strong-approx-eichler}
Let $N, n \geq 1$. Reduction induces a surjection from
$\GammaZero[lift]{N}$ onto the group of elements of determinant $1$ in
$\unit{(\EichlerOrder{N}/n\EichlerOrder{N})}$.
\end{lemma}

\begin{proof}
By strong approximation for $\SL{2}{}$, $\SL{2}{\QQ}$ is dense in
$\SL{2}{\FiniteAdeles{\QQ}}$; hence $\GammaZero[lift]{N} = \SL{2}{\QQ} \cap
\unit{\AdelicEichlerOrder{N}}$ is dense in the open subgroup of determinant
one elements of $\unit{\AdelicEichlerOrder{N}}$, which surjects onto the
determinant one elements of the finite quotient.
\end{proof}

\begin{proposition}[Intersections and products]\label{prop:lcm-gcd}
\slogan{$\GammaPrincipal{N}{M}$ satisfies the lcm and gcd laws in $M$.}
Let $N, M, M' \geq 1$ with $N$ squarefree. Then
\begin{align*}
  \GammaPrincipal[ad]{N}{M} \cap \GammaPrincipal[ad]{N}{M'} &=
  \GammaPrincipal[ad]{N}{\mathrm{lcm}(M,M')},\\
  \GammaPrincipal[ad]{N}{M}\,\GammaPrincipal[ad]{N}{M'} &=
  \GammaPrincipal[ad]{N}{\gcd(M,M')},
\end{align*}
and the same holds for the lifted versions.
\end{proposition}

\begin{proof}
Adelically the groups are images of products of local groups
$1 + p^{k}(\EichlerOrder{N}\otimes\pAdicIntegers{p})$, which form a decreasing
chain in $k$; intersections and products are computed prime by prime. For
the lifted versions, the intersection is immediate. For the product, both
groups contain $C := \GammaPrincipal[lift]{N}{\mathrm{lcm}(M,M')}$, which is
normal in $\GammaZero[lift]{N}$, and $\GammaZero[lift]{N}/C$ is the product
over $p$ of the determinant one parts of the local quotients by strong
approximation for Eichler orders (\cref{lem:strong-approx-eichler}); the
images of $\GammaPrincipal[lift]{N}{M}$ and $\GammaPrincipal[lift]{N}{M'}$ are
products of local subgroups, and the adelic argument applies.
\end{proof}

\subsection{The centred level}

\begin{definition}[Centred level]\label{def:centred-level}
Let $U \subseteq \AdelicGroup$ be a centred compact open subgroup with
$N := \HellingLevel{U}$. Let $\CentredLevelQuotient{U}$ be the least
$M \geq 1$ with $\GammaPrincipal[ad]{N}{M} \subseteq U$, and
$\CentredLevel{U} := N\cdot\CentredLevelQuotient{U}$ the \emph{centred
level}. For a centred $\Gamma \subseteq \RationalProjGroup$ commensurable
with $\PSL{2}{\ZZ}$ and containing a principal congruence subgroup, define
$\CentredLevelQuotient{\Gamma}$ and $\CentredLevel{\Gamma}$ in the same way
using $\GammaPrincipal[lift]{N}{M} \subseteq \lifted{\Gamma}$.
\end{definition}

\begin{lemma}[The centred level is well defined]\label{lem:centred-level-well-defined}
In the situation of \cref{def:centred-level}, the set of $M$ with
$\GammaPrincipal[ad]{N}{M} \subseteq U$ is nonempty and consists of the
multiples of $\CentredLevelQuotient{U}$; likewise in the classical case.
Moreover $\CentredLevel{U}$ is the level of the core
$\GammaPrincipal{N}{\CentredLevelQuotient{U}}$, i.e.\ the least $L$ with
$\GammaFull[lift]{L} \subseteq \GammaPrincipal[lift]{N}{\CentredLevelQuotient{U}}$,
and it only depends on the conjugacy class of $U$.
\end{lemma}

\begin{proof}
Nonempty since $U$ is open; closed under gcd by the product law
(\cref{prop:lcm-gcd}); closed under multiples since the groups decrease.
The statement about levels follows from the description of
$M\EichlerOrder{N}$ in \cref{lem:principal-lifted}. Conjugacy invariance:
centred conjugates are conjugate under $\MaxCompact{N}$ (\cref{prop:centring}),
which normalises $\GammaPrincipal[ad]{N}{M}$ (\cref{prop:principal-normal}).
\end{proof}

\begin{example}[Centred levels]\label{ex:centred-levels}
\begin{enumerate}
\item $\GammaZeroStar[ad]{N}$ and $\GammaZero[ad]{N}$ ($N$ squarefree) have
  Helling level $N$ and $\CentredLevelQuotient{} = 1$, so centred level $N$.
\item For $p \geq 3$ the centred conjugate of $\GammaZero[ad]{p^{2}}$ in
  \cref{ex:centring-examples} contains $\GammaFull[ad]{p}$ but not
  $\GammaPrincipal[ad]{1}{1}$, so it has centred level $p$.
\item For $p \geq 3$ the centred conjugate of $\GammaZero[ad]{p^{3}}$ in
  \cref{ex:centring-examples} contains $\GammaPrincipal[ad]{p}{p} =
  \{g \equiv 1 \bmod p,\ c \equiv 0 \bmod p^{2}\}$ but not
  $\GammaPrincipal[ad]{p}{1} = \GammaZero[ad]{p}$, so it has Helling level
  $p$ and centred level $p^{2}$, while the classical level of
  $\GammaZero{p^{3}}$ is $p^{3}$.
\end{enumerate}
\end{example}

\begin{remark}
If $U \subseteq \AdelicGroup$ is centred, then $\lifted{U}$ is centred with
the same Helling level and $\CentredLevelQuotient{\lifted{U}}$ divides
$\CentredLevelQuotient{U}$, since $\GammaPrincipal[ad]{N}{M} \subseteq U$
implies $\GammaPrincipal[lift]{N}{M} \subseteq \lifted{U}$.
\todo[inline]{The first claim needs the fixed set of the closure of
$\lifted{U}$ to have the same centre as $\FixedSet{U}$; check.}
\end{remark}
